% Generated draft: requires mathematical and source review.
\documentclass[11pt]{article}
\usepackage[margin=25mm]{geometry}
\usepackage{fontspec}
\setmainfont{Latin Modern Roman}
\usepackage{amsmath,amssymb,mathtools,unicode-math}
\setmathfont{Latin Modern Math}
\usepackage{enumitem,needspace,xurl,hyperref}
\setlength{\parindent}{0pt}
\setlength{\parskip}{5pt}
\setlength{\emergencystretch}{3em}
\newcommand{\N}{\mathbb{N}}
\newcommand{\Z}{\mathbb{Z}}
\newcommand{\Q}{\mathbb{Q}}
\newcommand{\R}{\mathbb{R}}
\newcommand{\C}{\mathbb{C}}
\newcommand{\F}{\mathbb{F}}
\newcommand{\A}{\mathbb{A}}
\newcommand{\PP}{\mathbb{P}}
\newcommand{\E}{\mathbb{E}}
\newcommand{\eps}{\varepsilon}
\newcommand{\dd}{\,\mathrm{d}}
\begin{document}
\section*{UnsolvedMath ID: 30003323}
\textbf{Source problem code:} catalogue entry
\subsection{Definitions and mathematical statement}
An ordered exponential field is an ordered field \(K\) equipped with an order-preserving group isomorphism \(\exp_K:(K,+)\to(K_{>0},\cdot)\). Work in the language \(\mathcal L_{\exp}=\{0,1,+,-,\cdot,<,\exp\}\). Such a structure is \emph{o-minimal} if every subset of \(K\) definable with parameters in one variable is a finite union of points and intervals.

The supplied record does not define the additional condition \(\mathrm{EXP}\), nor does it specify the valuation or the construction intended by ``residue exponential field.'' Consequently, the following is a statement template whose missing source conventions require verification.

Under the usual interpretation using the natural valuation, put
\[
\mathcal O_K=\{x\in K: |x|\le n\text{ for some }n\in\N_{>0}\},\qquad
\mathfrak m_K=\{x\in K: |x|<1/n\text{ for every }n\in\N_{>0}\},
\]
and let \(k_K=\mathcal O_K/\mathfrak m_K\), with its induced ordered-field structure. An induced residue exponential would be given by
\[
\overline{\exp}_K(x+\mathfrak m_K)=\exp_K(x)+\mathfrak m_K\qquad(x\in\mathcal O_K).
\]
This construction requires that the formula take values in the positive nonzero residue classes and be independent of the representative, and that it yield an ordered exponential field. Whether these requirements follow from the intended meaning of \(\mathrm{EXP}\) must be checked against the source; they are not additional hypotheses supplied by the record.

\emph{Main question in the supplied corrected statement.} For every o-minimal ordered exponential field \((K,\exp_K)\) satisfying the source condition \(\mathrm{EXP}\), is its intended residue exponential field \((k_K,\overline{\exp}_K)\) o-minimal?

\emph{Additional question retained from the original statement.} For every such \((K,\exp_K)\), does its residue exponential field also satisfy
\[
(k_K,\overline{\exp}_K)\equiv_{\mathcal L_{\exp}}(\R,\exp),
\]
where \(\exp\) on \(\R\) is the usual real exponential function and elementary equivalence means satisfaction of exactly the same first-order \(\mathcal L_{\exp}\)-sentences? The corrected statement retains only the o-minimality question; the record does not explain the omission of this additional target.
\subsection{Short English statement}
The main question asks whether o-minimality passes from an ordered exponential field satisfying the source condition \(\mathrm{EXP}\) to its residue exponential field. The original statement additionally asks whether that residue structure has the same first-order theory as the real exponential field. The meaning of \(\mathrm{EXP}\) and the intended residue construction remain unspecified.
\subsection{Sources}
\begin{itemize}
\item Supplied local UnsolvedMath registry record, permanent ID \(30003323\), title ``Residue Fields of O-Minimal Exponential Fields,'' source problem code \texttt{catalogue entry}, in the Oberwolfach Reports Open Problems collection. No individual proposer is supplied; the recorded proposed year is 2016.
\item Supplied source citation: \emph{Mini-Workshop: Surreal Numbers, Surreal Analysis, Hahn Fields and Derivations} (2017), Oberwolfach Reports, pp.~3313--3372, DOI \texttt{10.4171/owr/2016/60}. \url{https://doi.org/10.4171/owr/2016/60}. The record attributes the entry to workshop open-problem discussions.
\item Supplied literature source labeled ``Krapp'': \url{https://doi.org/10.4115/jla.2019.11.1}.
\item Supplied literature source labeled ``original OWR report'': \url{https://doi.org/10.4171/owr/2016/60}.
\item Supplied literature source labeled ``Krapp thesis overview'': \url{https://doi.org/10.1017/bsl.2021.64}.
\item Literature assessment dated exactly \texttt{2026-08-21}: the supplied curation labels the problem open, reports substantial residue-field and value-group results, and describes the extracted entry as a research overview whose underlying transfer questions remain open. This is supplied dated curation, not an independently verified literature assessment.
\item The registry labels the statement \texttt{corrected\_verified} and reports that the official passage was checked and rewritten to preserve its mathematical content. This is a supplied verification claim; the report passage and cited literature have not been independently consulted for this draft.
\end{itemize}
\end{document}
